
The study addressed three research questions. RQ1 (absolute performance): does logistic regression achieve $\geq 75$\% accuracy on held-out test? RQ2 (complexity value): how much does gradient boosting improve over logistic regression? RQ3 (feature importance): which metadata features predict maintenance status?

H-E1 (existence hypothesis): logistic regression trained on log-scaled GitHub metadata achieves $\geq 75$\% accuracy because repository maintenance status is linearly separable in transformed feature space. Pass criterion: accuracy $\geq 75$\% AND F1 $\geq 0.73$.

H-M1 (mechanism hypothesis): repository maintenance patterns form linearly separable clusters because maintained repositories exhibit consistently higher activity and lower staleness. Pass criteria: coefficient signs correct, performance gap $\leq 5$\%, feature overlap $\geq 2/3$.

Experimental protocol: (1) stratified 80/20 train/test split with random seed 42, (2) StandardScaler fitted on training set and applied to test set, (3) logistic regression trained with balanced class weights, (4) gradient boosting trained on identical split, (5) coefficient and feature importance extraction, (6) performance comparison.

Initial experiments included eight features: the six core features plus commit\_frequency\_median\_weekly and issue\_resolution\_rate. These derived features were later identified as tautological because they were constructed from formulas that encoded the binary label. Subsequent experiments restricted to six independently measured GitHub features.

Baseline comparison: a majority-class classifier (always predict ''maintained'') achieves 82.5\% accuracy by correctly predicting all 20 maintained test samples and missing all 4 abandoned samples. This establishes the trivial baseline that learned classifiers must exceed.
